% BANKS-SOURCE: pdf=87-99 print=77-89 kind=solution-section id=chapter07-implicit-solutions
\begin{bankssolutions}{Solutions to Implicit Exercises for Chapter 7}

% BANKS-SOURCE: pdf=87 print=77 kind=implicit-solution id=I-CH07-001
\BanksImplicitSolution{I-CH07-001}
\textit{(a) Finite functional identity.}
Write the source and field multiplets as columns.  For a finite symmetry
$R$, set $J'=R^{-T}J$ and make the change of integration variable
$\phi=R\psi$.  Invariance of the measure and action gives
\begin{align*}
 Z[J']
 &=\int[\dd\psi]\,
 \exp\left\{\ii S[\psi]
 +\ii\int\dd^D x\,(J')^TR\psi\right\}\\
 &=\int[\dd\psi]\,
 \exp\left\{\ii S[\psi]
 +\ii\int\dd^D x\,J^T\psi\right\}=Z[J].
\end{align*}
The branch of the logarithm connected continuously to the identity then
satisfies
\[
 W[R^{-T}J]=W[J].
\]

\textit{(b) Mean field and Legendre transform.}
The mean field transforms covariantly.  Indeed, regard
$J=R^TJ'$ while differentiating the preceding equality:
\[
 \varphi[J']
 =\frac{\delta W[J']}{\delta J'}
 =R\frac{\delta W[J]}{\delta J}
 =R\varphi[J].
\]
Banks's Legendre convention implies
$\delta\Gamma/\delta\varphi=-J$.  The source pairing is invariant,
$(J')^T(R\varphi)=J^T\varphi$, and hence
\begin{align*}
 \Gamma[R\varphi]
 &=W[J']-\int\dd^D x\,(J')^TR\varphi\\
 &=W[J]-\int\dd^D x\,J^T\varphi
 =\Gamma[\varphi].
\end{align*}

\textit{(c) Infinitesimal identities and possible failures.}
For $R=1+\epsilon T+O(\epsilon^2)$, the source variation is
$\delta J=-\epsilon T^TJ$.  The connected Ward identity is therefore
\[
 0=\int\dd^D x\,
 (T^TJ)_i(x)\frac{\delta W}{\delta J_i(x)}
 =\int\dd^D x\,J^T(x)T\varphi(x).
\]
The same identity in 1PI variables is
\[
 0=\int\dd^D x\,
 \frac{\delta\Gamma}{\delta\varphi_i(x)}(T\varphi)_i(x).
\]
For a symmetry-invariant zero-source state, differentiation shows that the
connected and 1PI vertices are invariant tensors.  When an infinite-volume
zero-source limit selects a broken vacuum, the same identity relates vertices
along the group orbit of that vacuum.

The derivation uses two hypotheses that matter in Banks's later discussion.
A nontrivial Jacobian of the measure contributes its local anomaly to the
right-hand side of the Ward identity.  A surviving integral of the current
divergence contributes a boundary term.  Either contribution changes the
functional identity and can obstruct the corresponding symmetry of $W$ and
$\Gamma$.

% BANKS-SOURCE: pdf=88 print=78 kind=implicit-solution id=I-CH07-002
\BanksImplicitSolution{I-CH07-002}
\textit{(a) Pole data.}
An invariant two-index tensor in the defining representation of $O(n)$ is
proportional to $\delta^{ab}$.  Consequently
\[
 G^{ab}(x-y)
 =\langle0|T\phi^a(x)\phi^b(y)|0\rangle
 =\delta^{ab}G(x-y).
\]
Its Lehmann representation may be written
\[
 \widetilde G^{ab}(p)=\delta^{ab}
 \int_0^\infty\dd\mu^2\,
 -\frac{\ii\rho(\mu^2)}{p^2+\mu^2-\ii\epsilon}.
\]
If the stable one-particle contribution is
$\rho(\mu^2)=Z_\phi\delta(\mu^2-m^2)+\rho_{\rm cont}(\mu^2)$, the same
$m$ and $Z_\phi$ occur for every component.  Phases of the particle states
can be chosen so that
\[
 \langle0|\phi^a(0)|p,b\rangle=Z_\phi^{1/2}\delta^{ab}.
\]
Thus the fields create an $n$-fold degenerate multiplet, and LSZ amputates
each external leg with the same pole residue.

\textit{(b) Tensor structure and crossing.}
The three invariant rank-four tensors of the full $O(n)$ group give
\[
 \mathcal M_{ab;cd}
 =\delta_{ab}\delta_{cd}F_s(s,t,u)
 +\delta_{ac}\delta_{bd}F_t(s,t,u)
 +\delta_{ad}\delta_{bc}F_u(s,t,u).
\]
All particles belong to a real representation, so every external leg is its
own crossed antiparticle.  Permuting legs in the amputated four-point
function relates the coefficient functions.  On defining
$A(s,t,u)=F_s(s,t,u)$, crossing gives
\begin{equation*}
 \boxed{\mathcal M_{ab;cd}(s,t,u)
 =\delta_{ab}\delta_{cd}A(s,t,u)
 +\delta_{ac}\delta_{bd}A(t,s,u)
 +\delta_{ad}\delta_{bc}A(u,t,s).}
\end{equation*}
Bose symmetry for either incoming or outgoing pair requires
\[
 A(s,t,u)=A(s,u,t).
\]
For example, the complete exchange and crossing statements include
\begin{align*}
 \mathcal M_{ab;cd}(s,t,u)
 &=\mathcal M_{ba;cd}(s,u,t)
 =\mathcal M_{ab;dc}(s,u,t),\\
 \mathcal M_{ab;cd}(s,t,u)
 &=\mathcal M_{ac;bd}(t,s,u)
 =\mathcal M_{ad;cb}(u,t,s),
\end{align*}
where the last line uses analytic continuation between the physical channel
regions.

\textit{(c) Irreducible channels.}
The $s$-channel projectors are
\begin{align*}
 (P_{\mathbf1})_{ab;cd}&=\frac1n\delta_{ab}\delta_{cd},\\
 (P_{\mathrm A})_{ab;cd}
 &=\frac12(\delta_{ac}\delta_{bd}-\delta_{ad}\delta_{bc}),\\
 (P_{\mathrm S})_{ab;cd}
 &=\frac12(\delta_{ac}\delta_{bd}+\delta_{ad}\delta_{bc})
 -\frac1n\delta_{ab}\delta_{cd}.
\end{align*}
Writing
$\mathcal M=T_{\mathbf1}P_{\mathbf1}+T_{\mathrm A}P_{\mathrm A}
+T_{\mathrm S}P_{\mathrm S}$ and abbreviating
$A_s=A(s,t,u)$, $A_t=A(t,s,u)$, $A_u=A(u,t,s)$ gives
\[
 T_{\mathbf1}=nA_s+A_t+A_u,
 \qquad T_{\mathrm A}=A_t-A_u,
 \qquad T_{\mathrm S}=A_t+A_u.
\]
These expressions supply a check: an exchange of the two incoming bosons is
even in the singlet and symmetric-traceless channels and odd in the
antisymmetric channel, together with $t\leftrightarrow u$.

\textit{(d) Replacing $O(n)$ by $SO(n)$.}
The Levi-Civita tensor changes sign under an $O(n)$ reflection, so a single
epsilon tensor is absent from the preceding basis.  For $SO(4)$,
$\epsilon_{abcd}E(s,t,u)$ is an additional invariant structure, with its
coefficient constrained by the same leg permutations and the signs of the
epsilon tensor.  For $SO(2)$, tensors of the form
$\epsilon_{ab}\delta_{cd}$ and their index permutations add parity-odd
rank-four structures, subject to the two-dimensional identities.  A single
epsilon has the wrong rank for $n=3$ and for $n>4$, while contractions of two
epsilon tensors reduce to Kronecker deltas.  Thus there is no further
rank-four structure in those dimensions.

% BANKS-SOURCE: pdf=88-91 print=78-81 kind=implicit-solution id=I-CH07-003
\BanksImplicitSolution{I-CH07-003}
We treat the eight statements in order.

\begin{enumerate}[label=(\alph*)]
\item Let
$\mathcal O=\phi_{i_1}(x_1)\cdots\phi_{i_N}(x_N)$.  A local translation
$\delta\phi_i=-f^\nu\partial_\nu\phi_i$ changes the action by the defining
Noether formula
\[
 \delta S=\int\dd^D x\,\partial_\mu f_\nu(x)T^{\mu\nu}(x).
\]
Changing variables in the path integral gives, to first order in $f$,
\[
 0=\langle\delta\mathcal O\rangle
 +\ii\langle T(\mathcal O\,\delta S)\rangle.
\]
The variation of the inserted fields is
\[
 \langle\delta\mathcal O\rangle
 =-\sum_{k=1}^N f_\nu(x_k)
 \left\langle T\left(
 \phi_{i_1}(x_1)\cdots
 \partial^\nu\phi_{i_k}(x_k)\cdots
 \phi_{i_N}(x_N)\right)\right\rangle.
\]
After integrating the action variation by parts, the coefficient of the
arbitrary function $f_\nu(x)$ must vanish.  This yields
\begin{align*}
 &\partial_\mu^x
 \left\langle0\left|T\left(
 \begin{aligned}
  &T^{\mu\nu}(x)\phi_{i_1}(x_1)\cdots\\[-2pt]
  &{}\qquad\phi_{i_N}(x_N)
 \end{aligned}
 \right)\right|0\right\rangle\\[-2pt]
 ={}&+\ii\sum_{k=1}^N\delta^D(x-x_k)\\
 &\mathrel{\phantom{=}}\times\\[-2pt]
 &\qquad
 \left\langle0\left|T\left(
 \begin{aligned}
  &\phi_{i_1}(x_1)\cdots
   \partial^\nu\phi_{i_k}(x_k)\cdots\\[-2pt]
  &{}\qquad\phi_{i_N}(x_N)
 \end{aligned}
 \right)\right|0\right\rangle.
\end{align*}
The delta functions are the contact terms generated when the translation
acts on each insertion.  At separated points they disappear and the stress
tensor insertion is conserved.  A Jacobian for the change of variables would
add the corresponding anomaly insertion to this identity.

\item Put
$T'^{\mu\nu}=T^{\mu\nu}+\partial_\lambda
M^{\mu\nu\lambda}$.  Commuting derivatives and using the antisymmetry in
$\mu$ and $\lambda$ gives
\[
 \partial_\mu T'^{\mu\nu}
 =\partial_\mu T^{\mu\nu}
 +\frac12\partial_\mu\partial_\lambda
 \left(M^{\mu\nu\lambda}+M^{\lambda\nu\mu}\right)
 =\partial_\mu T^{\mu\nu}.
\]
For the charge, $M^{0\nu0}=0$, and therefore
\begin{align*}
 P'^{\nu}-P^\nu
 &=\int\dd^{D-1}\mathbf x\,\partial_\lambda M^{0\nu\lambda}
 =\int\dd^{D-1}\mathbf x\,\partial_iM^{0\nu i}\\
 &=\lim_{R\to\infty}\int_{S_R^{D-2}}\dd S_i\,M^{0\nu i}=0.
\end{align*}
The last equality is precisely the assumed spatial falloff.  This is the
Noether improvement ambiguity used throughout the section.

\item Consider an orientation-preserving change of coordinates $x\mapsto y$.
A scalar and a one-form transform according to
\[
 \phi'(y)=\phi(x),\qquad
 A'_\alpha(y)=\frac{\partial x^\mu}{\partial y^\alpha}A_\mu(x),
\]
while
\[
 g'_{\alpha\beta}(y)
 =\frac{\partial x^\mu}{\partial y^\alpha}
  \frac{\partial x^\nu}{\partial y^\beta}g_{\mu\nu}(x).
\]
Taking determinants gives
\[
 \sqrt{-g'(y)}
 =\left|\det\frac{\partial x}{\partial y}\right|\sqrt{-g(x)},
 \qquad
 \dd^D y\sqrt{-g'(y)}=\dd^D x\sqrt{-g(x)}.
\]
The inverse metric carries the inverse Jacobians, so
\[
 g'^{\alpha\beta}\partial'_{\alpha}\phi'
 \partial'_{\beta}\phi'
 =g^{\mu\nu}\partial_\mu\phi\partial_\nu\phi.
\]
The potential is a scalar as well.  Hence every factor in $S_\phi$ combines
to make an invariant density.

The exterior derivative of the one-form is a two-form:
\[
 F'_{\alpha\beta}(y)
 =\frac{\partial x^\mu}{\partial y^\alpha}
  \frac{\partial x^\nu}{\partial y^\beta}F_{\mu\nu}(x).
\]
Its two Jacobian pairs cancel those of the two inverse metrics in
$g^{\mu\rho}g^{\nu\sigma}F_{\mu\nu}F_{\rho\sigma}$.  Combining this scalar
with the invariant volume form proves invariance of $S_M$.  The calculation
also explains why $\sqrt{-g}$ and two inverse metrics occur in the curved
Maxwell action.

\item An infinitesimal conformal coordinate transformation
$x^\mu\mapsto x^\mu+f^\mu(x)$ obeys
\begin{equation*}
 \partial_\mu f_\nu+\partial_\nu f_\mu
 =\frac12\eta_{\mu\nu}\partial_\rho f^\rho.
\end{equation*}
Define $\sigma=\frac14\partial\cdot f$.  Differentiate the conformal Killing
equation and combine three permutations of its indices to obtain
\[
 \partial_\mu\partial_\nu f_\rho
 =\eta_{\mu\rho}\partial_\nu\sigma
 +\eta_{\nu\rho}\partial_\mu\sigma
 -\eta_{\mu\nu}\partial_\rho\sigma.
\]
Further differentiation in dimension greater than two implies
$\partial_\mu\partial_\nu\sigma=0$.  Thus $\sigma$ is affine and $f$ is at
most quadratic.  Integrating the preceding equation gives the complete
solution
\begin{equation*}
 f^\mu(x)=a^\mu+\omega^\mu{}_{\nu}x^\nu+\lambda x^\mu
 +b^\mu x^2-2x^\mu(b\cdot x),
 \qquad \omega_{\mu\nu}=-\omega_{\nu\mu}.
\end{equation*}
The parameters count as $4+6+1+4=15$.  They generate translations, Lorentz
transformations, dilatations, and special conformal transformations.  For
the special conformal term,
$\partial\cdot f=-8b\cdot x$, which checks the conformal Killing equation.

Use the following differential generators, with the sign of $K_\mu$ chosen
to match Banks's special conformal transformation:
\[
 P_\mu=\partial_\mu,\quad
 M_{\mu\nu}=x_\mu\partial_\nu-x_\nu\partial_\mu,\quad
 D=x\cdot\partial,\quad
 K_\mu=x^2\partial_\mu-2x_\mu x\cdot\partial.
\]
Direct evaluation of the vector-field commutators gives
\begin{align*}
 [M_{\mu\nu},M_{\rho\sigma}]
 & =\eta_{\nu\rho}M_{\mu\sigma}
 -\eta_{\mu\rho}M_{\nu\sigma}
 -\eta_{\nu\sigma}M_{\mu\rho}
 +\eta_{\mu\sigma}M_{\nu\rho},\\
 [M_{\mu\nu},P_\rho]
 &=\eta_{\nu\rho}P_\mu-\eta_{\mu\rho}P_\nu,\\
 [M_{\mu\nu},K_\rho]
 &=\eta_{\nu\rho}K_\mu-\eta_{\mu\rho}K_\nu,\\
 [D,P_\mu]&=-P_\mu,\qquad [D,K_\mu]=K_\mu,\\
 [P_\mu,K_\nu]&=2M_{\mu\nu}-2\eta_{\mu\nu}D,\\
 [P_\mu,P_\nu]&=[K_\mu,K_\nu]=[D,M_{\mu\nu}]=0.
\end{align*}
To identify the algebra, introduce $J_{AB}=-J_{BA}$ in six dimensions with
metric
$\widehat\eta_{AB}=\operatorname{diag}(\eta_{\mu\nu},+1,-1)$ and brackets
\[
 [J_{AB},J_{CD}]
 =\widehat\eta_{BC}J_{AD}-\widehat\eta_{AC}J_{BD}
 -\widehat\eta_{BD}J_{AC}+\widehat\eta_{AD}J_{BC}.
\]
The assignments
\[
 M_{\mu\nu}=J_{\mu\nu},\qquad D=J_{45},\qquad
 P_\mu=J_{\mu4}+J_{\mu5},\qquad
 K_\mu=J_{\mu5}-J_{\mu4}
\]
reproduce the displayed conformal brackets.  The six-dimensional metric has
signature $(2,4)$, proving the isomorphism with $\mathfrak{so}(2,4)$.

For a symmetric conserved stress tensor, the current associated with a
conformal Killing vector is $j_f^\mu=T^{\mu\nu}f_\nu$.  The conformal Killing
equation gives
\[
 \partial_\mu j_f^\mu
 =T^{\mu\nu}\partial_\mu f_\nu
 =\frac14(\partial\cdot f)T^\mu{}_{\mu}.
\]
The dilation has $\partial\cdot f=4\lambda$, while a special conformal
transformation has a position-dependent divergence.  Their currents are
conserved when the improved stress tensor is traceless.  Equivalently, a
local Weyl variation $\delta g_{\mu\nu}=2\sigma(x)g_{\mu\nu}$ couples to
$T^\mu{}_{\mu}$; invariance for arbitrary $\sigma(x)$ sets this trace to
zero.  A removable virial divergence is absorbed into the same improvement.

\item Let $D$ be the Hermitian dilation charge.  With a conventional choice
of signs, its commutator with momentum is
\[
 [D,P_\mu]=\ii P_\mu.
\]
Exponentiation shows that a dilation sends a momentum eigenstate to another
one whose four-momentum is multiplied by a continuous positive factor.
Consequently
\[
 P^2|p\rangle=-m^2|p\rangle
 \quad\Longrightarrow\quad
 P^2U(\lambda)|p\rangle=-\lambda^2m^2U(\lambda)|p\rangle
\]
after a harmless choice between $\lambda$ and $\lambda^{-1}$ in the group
parameter.  A positive isolated mass would therefore generate a continuous
orbit of distinct positive masses.  This conflicts with a discrete
one-particle mass spectrum.  The value $m^2=0$ is fixed under every dilation,
so all particles in such a conformal theory are massless.  This argument
assumes that the conformal symmetry acts on the vacuum and on asymptotic
one-particle states.

\item Since $g_{\mu\nu}=g_{\nu\mu}$, its functional derivative is symmetric:
$T_{\mathrm H}^{\mu\nu}=T_{\mathrm H}^{\nu\mu}$.  Under an infinitesimal
diffeomorphism generated by $\xi^\mu$,
\[
 \delta_\xi g_{\mu\nu}=\nabla_\mu\xi_\nu+\nabla_\nu\xi_\mu.
\]
On the matter equations of motion, the matter-field variation contributes
zero.  Diffeomorphism invariance and integration by parts then give
\begin{align*}
 0=\delta_\xi S
 &=-\int\dd^D x\sqrt{-g}\,
 T_{\mathrm H}^{\mu\nu}\nabla_\mu\xi_\nu\\
 &=\int\dd^D x\sqrt{-g}\,
 (\nabla_\mu T_{\mathrm H}^{\mu\nu})\xi_\nu.
\end{align*}
Arbitrariness of $\xi_\nu$ proves
$\nabla_\mu T_{\mathrm H}^{\mu\nu}=0$.

For flat-space fields $\Phi_A$, define the canonical tensor and spin current
by
\[
 T_{\mathrm{can}}^{\mu\nu}
 =-\frac{\partial\mathcal L}{\partial(\partial_\mu\Phi_A)}
 \partial^\nu\Phi_A+\eta^{\mu\nu}\mathcal L,
 \qquad
 S^{\lambda\mu\nu}
 =-\frac{\partial\mathcal L}{\partial(\partial_\lambda\Phi_A)}
 (\Sigma^{\mu\nu})_A{}^B\Phi_B.
\]
Here $S^{\lambda\mu\nu}=-S^{\lambda\nu\mu}$.  Lorentz invariance gives
\[
 \partial_\lambda S^{\lambda\mu\nu}
 =T_{\mathrm{can}}^{\nu\mu}-T_{\mathrm{can}}^{\mu\nu}
\]
on shell.  Set
\[
 B^{\lambda\mu\nu}
 =\frac12\left(S^{\lambda\mu\nu}+S^{\mu\nu\lambda}
 +S^{\nu\mu\lambda}\right),
 \qquad B^{\lambda\mu\nu}=-B^{\mu\lambda\nu},
\]
and
\[
 T_{\mathrm B}^{\mu\nu}
 =T_{\mathrm{can}}^{\mu\nu}+\partial_\lambda
 B^{\lambda\mu\nu}.
\]
The antisymmetry of $B$ preserves conservation and the translation charges,
as in part (b).  The Lorentz identity makes $T_{\mathrm B}$ symmetric.  The
flat-space limit of the metric tensor equals this Belinfante tensor on shell,
up to terms proportional to field equations and further identically
conserved improvements.  This is the precise relation between the metric and
Noether constructions.

Using symmetry and conservation,
\begin{align*}
 \partial_\lambda J^{\lambda\mu\nu}
 &=T_{\mathrm B}^{\mu\nu}-T_{\mathrm B}^{\nu\mu}
 +x^\mu\partial_\lambda T_{\mathrm B}^{\lambda\nu}
 -x^\nu\partial_\lambda T_{\mathrm B}^{\lambda\mu}=0.
\end{align*}
Thus the spatial integrals of $J^{0\mu\nu}$ generate rotations and boosts.
For example, the boost charge
\[
 K^i=\int\dd^{D-1}\mathbf x\,J^{00i}
 =tP^i-\int\dd^{D-1}\mathbf x\,x^iT_{\mathrm B}^{00}
\]
has the explicit derivative $\partial K^i/\partial t=P^i$.  Conservation and
the Heisenberg equation
$\dd K^i/\dd t=\partial K^i/\partial t+\ii[H,K^i]=0$ therefore give
$[H,K^i]=\ii P^i$.  This is Banks's stated reason that a conserved boost
charge need not commute with the Hamiltonian.

Banks displays
$T_{\text{source}}^{\mu\nu}=(-g)^{-1/2}\delta S/\delta g_{\mu\nu}$.
Under the variational definition in the exercise,
$T_{\text{source}}^{\mu\nu}=T_{\mathrm H}^{\mu\nu}/2$.  The factor is fixed
before identifying the metric derivative with the normalized translation
current.  The symmetry and conservation arguments above use
$T_{\mathrm H}^{\mu\nu}$.

\item For the Hermitian massless Dirac Lagrangian
\[
 \mathcal L_D=-\frac{\ii}{2}\bar\psi\gamma^\rho
 \overleftrightarrow{\partial_\rho}\psi,
 \qquad
 \bar\psi\gamma^\rho\overleftrightarrow{\partial_\rho}\psi
 =\bar\psi\gamma^\rho\partial_\rho\psi
 -(\partial_\rho\bar\psi)\gamma^\rho\psi,
\]
the canonical tensor is
\[
 T_{D,\mathrm{can}}^{\mu\nu}
 =-\frac{\ii}{2}\bar\psi\gamma^\mu
 \overleftrightarrow{\partial^\nu}\psi
 -\eta^{\mu\nu}\mathcal L_D.
\]
With
$\sigma^{\mu\nu}=\frac{\ii}{2}[\gamma^\mu,\gamma^\nu]$, the spin current is
\[
 S_D^{\lambda\mu\nu}
 =-\frac14\bar\psi\{\gamma^\lambda,\sigma^{\mu\nu}\}\psi.
\]
Substitution in the Belinfante formula from part (f), followed by the Dirac
equations and $\mathcal L_D=0$ on shell, gives
\[
 T_D^{\mu\nu}=-\frac{\ii}{4}\bar\psi\left(
 \gamma^\mu\overleftrightarrow{\partial^\nu}
 +\gamma^\nu\overleftrightarrow{\partial^\mu}\right)\psi.
\]
Its symmetry is explicit.  Differentiation, followed by use of
$-\ii\gamma^\rho\partial_\rho\psi=0$ and
$(\partial_\rho\bar\psi)\ii\gamma^\rho=0$, proves
$\partial_\mu T_D^{\mu\nu}=0$.  Its trace is
\[
 T_{D\,\mu}{}^\mu
 =-\frac{\ii}{2}\bar\psi\gamma^\mu
 \overleftrightarrow{\partial_\mu}\psi=0
\]
on the same equations of motion.

For $\mathcal L_M=-\frac14F_{\rho\sigma}F^{\rho\sigma}$, the canonical tensor
is
\[
 T_{M,\mathrm{can}}^{\mu\nu}
 =+F^{\mu\lambda}\partial^\nu A_\lambda
 -\frac14\eta^{\mu\nu}F_{\rho\sigma}F^{\rho\sigma}.
\]
Add the improvement $-\partial_\lambda(F^{\mu\lambda}A^\nu)$.  The Maxwell
equation $\partial_\lambda F^{\mu\lambda}=0$ then gives
\[
 +F^{\mu\lambda}\partial^\nu A_\lambda
 -F^{\mu\lambda}\partial_\lambda A^\nu
 =+F^{\mu\lambda}F^\nu{}_{\lambda}.
\]
The resulting gauge-invariant Belinfante tensor is
\[
 T_M^{\mu\nu}
 =+F^{\mu\lambda}F^\nu{}_{\lambda}
 -\frac14\eta^{\mu\nu}F_{\rho\sigma}F^{\rho\sigma}.
\]
It is symmetric.  Its divergence is
\begin{align*}
 \partial_\mu T_M^{\mu\nu}
 ={}&+(\partial_\mu F^{\mu\lambda})F^\nu{}_{\lambda}
 +F^{\mu\lambda}\partial_\mu F^\nu{}_{\lambda}
 -\frac12F^{\rho\sigma}\partial^\nu F_{\rho\sigma}=0.
\end{align*}
The first term vanishes by $\partial_\mu F^{\mu\lambda}=0$; the remaining
terms cancel by the Bianchi identity.  Finally,
\[
 T_{M\,\mu}{}^\mu
 =+F^{\mu\lambda}F_{\mu\lambda}
 -\frac44F_{\rho\sigma}F^{\rho\sigma}=0.
\]
The trace cancellation is special to four dimensions, as expected from the
dimensionless Maxwell coupling.

\item The free scalar Belinfante tensor is already symmetric:
\[
 T_{\mathrm B}^{\mu\nu}
 =\partial^\mu\phi\partial^\nu\phi
 -\frac12\eta^{\mu\nu}(\partial\phi)^2.
\]
Its divergence is $(\partial^\nu\phi)\partial^2\phi$, and its trace in $D$
dimensions is
\[
 T_{\mathrm B\,\mu}{}^\mu=-\frac{D-2}{2}(\partial\phi)^2.
\]
The improvement is identically conserved because derivatives commute:
\[
 \partial_\mu(\partial^\mu\partial^\nu
 -\eta^{\mu\nu}\partial^2)\phi^2=0.
\]
Taking its trace and using $\partial^2\phi=0$ gives, with
$\xi_D=(D-2)/[4(D-1)]$,
\begin{align*}
 T_{\mathrm I\,\mu}{}^\mu
 &=-\frac{D-2}{2}(\partial\phi)^2
 +\xi_D(D-1)\partial^2\phi^2\\
 &=-\frac{D-2}{2}(\partial\phi)^2
 +\frac{D-2}{2}\left[(\partial\phi)^2+\phi\partial^2\phi\right]=0.
\end{align*}
Thus $T_{\mathrm I}$ is conserved and traceless on shell.

The canonical tensor depends only on derivatives of $\phi$ and is invariant
under $\phi\mapsto\phi+c$.  The improvement changes by
\begin{align*}
 \Delta_cT_{\mathrm I}^{\mu\nu}
 &=-\xi_D(\partial^\mu\partial^\nu
 -\eta^{\mu\nu}\partial^2)(2c\phi+c^2)\\
 &=-2c\xi_D(\partial^\mu\partial^\nu
 -\eta^{\mu\nu}\partial^2)\phi.
\end{align*}
This local tensor is generally nonzero, even on the massless field equation.
It displays the explicit loss of constant field-shift invariance emphasized
by Banks.
\end{enumerate}

% BANKS-SOURCE: pdf=97 print=87 kind=implicit-solution id=I-CH07-004
\BanksImplicitSolution{I-CH07-004}
\textit{(a) Chiral group, vacuum, and coset.}
The massless quark Lagrangian separates into
\[
 \mathcal L_q=-\ii\bar q_L\slashed Dq_L
 -\ii\bar q_R\slashed Dq_R.
\]
It is invariant under independent flavor rotations
$q_L\mapsto Lq_L$ and $q_R\mapsto Rq_R$.  The classical flavor group is
$U(2)_L\times U(2)_R$.  The axial singlet factor has a quantum anomaly, while
the vector singlet is baryon number and acts trivially on the pion field.
The relevant exact non-abelian group is therefore
\[
 G=SU(2)_L\times SU(2)_R.
\]
Choose the order-parameter orientation
$\Phi\sim\bar q_Lq_R$, with flavor indices arranged so that
$\Phi\mapsto L^\dagger\Phi R$.  A vacuum
$\langle\Phi\rangle=v\mathbf1$ is fixed when $L=R$, so
\[
 H=SU(2)_V,\qquad
 \frac{G}{H}=\frac{SU(2)_L\times SU(2)_R}{SU(2)_V}
 \simeq SU(2)\simeq S^3.
\]
Its dimension is three, giving the pion triplet.

A useful coset coordinate is an $SU(2)$ matrix $\Sigma$.  In Banks's passive
flavor convention,
\[
 \Sigma\longmapsto L^\dagger\Sigma R,\qquad
 \Sigma=\exp\left(\frac{\ii\Pi}{f}\right),\qquad
 \Pi=\pi^at^a.
\]
The action is transitive and the stabilizer of $\Sigma=\mathbf1$ has $L=R$,
which supplies the coset identification above.  Replacing $\Sigma$ by
$\Sigma^\dagger$ exchanges this convention with the common active one.

\textit{(b) Leading Lagrangian and pion expansion.}
The chiral group acts transitively on $G/H$, so every invariant function
without derivatives is constant.  At two derivatives, the invariant
bilinear form on the simple algebra $\mathfrak{su}(2)$ is unique up to scale.
Parity exchanges left and right.  These observations give
\[
 \mathcal L_2=-f^2\Tr(\partial_\mu\Sigma^\dagger
 \partial^\mu\Sigma).
\]

To expand it, use the Maurer--Cartan series
\begin{align*}
 \Sigma^\dagger\partial_\mu\Sigma
 ={}&\frac{\ii}{f}\partial_\mu\Pi
 +\frac{1}{2f^2}[\Pi,\partial_\mu\Pi]
 -\frac{\ii}{6f^3}[\Pi,[\Pi,\partial_\mu\Pi]]+O(\pi^4).
\end{align*}
Since
$\Tr(\partial_\mu\Sigma^\dagger\partial^\mu\Sigma)
=-\Tr[(\Sigma^\dagger\partial_\mu\Sigma)^2]$, and
$[t^a,t^b]=\ii\epsilon^{abc}t^c$, one obtains
\begin{equation*}
 \boxed{\mathcal L_2
 =-\frac12\partial_\mu\pi^a\partial^\mu\pi^a
 -\frac{1}{24f^2}\left[
 (\pi^a\partial_\mu\pi^a)^2
 -(\pi^a\pi^a)(\partial_\mu\pi^b\partial^\mu\pi^b)
 \right]+O(\pi^6).}
\end{equation*}
The quadratic term confirms canonical normalization.  Odd powers vanish by
the $\pi\mapsto-\pi$ parity of the pion-only action.

\textit{(c) Currents and the four-pion amplitude.}
Promote $L$ and $R$ separately to local transformations while retaining only
terms proportional to derivatives of their parameters.  For
$L=\exp(\ii\alpha_L^at^a)$ and
$R=\exp(\ii\alpha_R^at^a)$, the variation takes the form
$\delta\mathcal L_2=(\partial_\mu\alpha_L^a)J_L^{a\mu}
+(\partial_\mu\alpha_R^a)J_R^{a\mu}$, with
\begin{align*}
 J_L^{a\mu}&=-2\ii f^2\Tr\left(t^a
 \partial^\mu\Sigma\,\Sigma^\dagger\right),\\
 J_R^{a\mu}&=+2\ii f^2\Tr\left(t^a
 \Sigma^\dagger\partial^\mu\Sigma\right).
\end{align*}
The field equation makes both currents conserved.  Expanding gives
\begin{align*}
 J_L^{a\mu}
 &=+f\partial^\mu\pi^a
 -\frac12\epsilon^{abc}\pi^b\partial^\mu\pi^c+O(\pi^3),\\
 J_R^{a\mu}
 &=-f\partial^\mu\pi^a
 -\frac12\epsilon^{abc}\pi^b\partial^\mu\pi^c+O(\pi^3).
\end{align*}
With the local-parameter signs chosen above, the vector and axial
currents are
\begin{align*}
 V^{a\mu}=J_L^{a\mu}+J_R^{a\mu}
 &=-\epsilon^{abc}\pi^b\partial^\mu\pi^c+O(\pi^4),\\
 A^{a\mu}=J_L^{a\mu}-J_R^{a\mu}
 &=+2f\partial^\mu\pi^a+O(\pi^3).
\end{align*}
Changing the signs assigned to the infinitesimal generators changes all
current signs together and leaves their Ward identities unchanged.

The quartic term produces the on-shell tree amplitude
\[
 s=-(p_1+p_2)^2,\qquad t=-(p_1-p_3)^2,\qquad
 u=-(p_1-p_4)^2,\qquad s+t+u=0.
\]
\begin{equation*}
 \boxed{\mathcal M_{ab;cd}(s,t,u)
 =\frac{1}{4f^2}\left(
 \delta_{ab}\delta_{cd}s
 +\delta_{ac}\delta_{bd}t
 +\delta_{ad}\delta_{bc}u\right),}
\end{equation*}
This has exactly the $O(3)$ tensor and crossing form found in
Exercise~\ref{exercise:I-CH07-002}.  Taking any external pion momentum to
zero sends $s,t,u$ to zero and therefore makes the amplitude vanish.  This is
the Adler zero.  It follows here from the derivative couplings forced by the
Goldstone symmetry.

\textit{(d) Current normalization, weak decay, and the soft theorem.}
For a canonically normalized one-pion state, the linear current terms imply
\[
 \langle0|J_L^{a\mu}(0)|\pi^b(p)\rangle
 =\ii f p^\mu\delta^{ab},
 \qquad
 \langle0|A^{a\mu}(0)|\pi^b(p)\rangle
 =\ii F_\pi p^\mu\delta^{ab},
 \qquad F_\pi=2f,
\]
where a pion-state phase fixes the displayed signs.  Banks's parametrization
uses $f$ in the exponential and in the left-current matrix element.  A
convention that defines the decay constant with the full axial current uses
$F_\pi=2f$.  Equivalently, writing the standard chiral Lagrangian as
$-F_\pi^2\Tr(\partial\Sigma^\dagger\partial\Sigma)/4$ and its exponential as
$\exp(\ii\tau^a\pi^a/F_\pi)$ gives the same theory.

Coupling the charged left current to the weak interaction and integrating
out the $W$ gives, in the standard full-current normalization,
\begin{align*}
 \mathcal M(\pi^-\to\ell^-\bar\nu_\ell)
 &=\frac{G_FV_{ud}}{\sqrt2}F_\pi p_\mu
 \bar u_\ell\gamma^\mu(1-\gamma_5)v_\nu\\
 &=\frac{G_FV_{ud}}{\sqrt2}F_\pi m_\ell
 \bar u_\ell(1-\gamma_5)v_\nu.
\end{align*}
The second line follows from $p=p_\ell+p_\nu$ and the lepton Dirac
equations.  It exhibits helicity suppression.  With small quark masses
restored so that the physical pion can decay, the leading width is
\[
 \Gamma(\pi^-\to\ell^-\bar\nu_\ell)
 =\frac{G_F^2|V_{ud}|^2F_\pi^2m_\ell^2m_\pi}{8\pi}
 \left(1-\frac{m_\ell^2}{m_\pi^2}\right)^2,
\]
before electromagnetic and higher-order chiral corrections.

The same matrix element fixes the pole in any Green function containing the
axial current:
\[
 \widetilde G_{A\mathcal O}^{a\mu}(q)
 =-\frac{\ii F_\pi q^\mu}{q^2-\ii\epsilon}
 \mathcal M(\pi^a(q),\mathcal O)+\text{terms regular at }q^2=0.
\]
Insert this expression into the axial Ward identity and take $q\to0$.  With
the same phase convention, the result is the current-algebra relation
\[
 \lim_{q\to0}\langle\beta;\pi^a(q)|\mathcal O|\alpha\rangle
 =-\frac{\ii}{F_\pi}
 \langle\beta|[Q_5^a,\mathcal O]|\alpha\rangle
 +\text{external-particle pole terms}.
\]
It relates soft-pion emission from other hadrons to the chiral transformation
of the amplitude without the soft pion.  For a regular pion-only amplitude,
the commutator structure and derivative action reproduce the Adler zero
found directly above.

% BANKS-SOURCE: pdf=99 print=89 kind=implicit-solution id=I-CH07-005
\BanksImplicitSolution{I-CH07-005}
\textit{(a) Finite connection law.}
Introduce the Maurer--Cartan form
\[
 \Theta_\mu=g^{-1}\partial_\mu g.
\]
For $g'=gh$, the product rule gives
\begin{align*}
 \Theta'_\mu
 &=(gh)^{-1}\partial_\mu(gh)\\
 &=h^{-1}g^{-1}(\partial_\mu g)h+h^{-1}\partial_\mu h\\
 &=h^{-1}\Theta_\mu h+h^{-1}\partial_\mu h.
\end{align*}
The second term belongs to $\mathfrak h$.  The assumed
$\operatorname{Ad}(H)$-invariance of
$\mathfrak g=\mathfrak h\oplus\mathfrak m$ means that projection commutes
with conjugation by $h$:
\[
 P_{\mathfrak h}(h^{-1}Xh)=h^{-1}P_{\mathfrak h}(X)h.
\]
Projecting the transformed Maurer--Cartan form therefore gives
\begin{equation*}
 \boxed{A'_\mu=h^{-1}A_\mu h+h^{-1}\partial_\mu h
 =h^{-1}(A_\mu+\partial_\mu)h.}
\end{equation*}
This is the inhomogeneous transformation law of a connection for the local
right $H$ action.

\textit{(b) Components and covariant checks.}
Choose the basis in the exercise anti-Hermitian and write
$[t_n,t_p]=f_{np}{}^m t_m$ and
$h=\exp(\epsilon^mt_m)$.  To first order,
\[
 \delta A_\mu=\partial_\mu\epsilon+[A_\mu,\epsilon],
 \qquad
 \delta A_\mu^m=\partial_\mu\epsilon^m
 +f_{np}{}^mA_\mu^n\epsilon^p.
\]
Hermitian generators move factors of $\ii$ between the structure constants,
the group parameter, and the connection; the matrix transformation above
fixes those convention-dependent factors unambiguously.

The covariant derivative supplies a direct check:
\begin{align*}
 (D_\mu g)'
 &=\partial_\mu(gh)-(gh)A'_\mu\\
 &=(\partial_\mu g)h+g\partial_\mu h
 -gh\left(h^{-1}A_\mu h+h^{-1}\partial_\mu h\right)\\
 &=(\partial_\mu g-gA_\mu)h=(D_\mu g)h.
\end{align*}
Substitution of the finite connection law shows that the derivative terms in
$h$ cancel in the curvature
$F_{\mu\nu}=\partial_\mu A_\nu-\partial_\nu A_\mu+[A_\mu,A_\nu]$
and leaves $F_{\mu\nu}'=h^{-1}F_{\mu\nu}h$.

\textit{(c) A general generator metric.}
For a general generator metric
$\kappa_{mn}=\Tr(t_mt_n)$, the projection coefficients are
\[
 A_\mu^m=(\kappa^{-1})^{mn}
 \Tr(t_n g^{-1}\partial_\mu g).
\]
In the common normalization $\kappa_{mn}=\kappa\delta_{mn}$, Banks's
component formula acquires the factor $1/\kappa$.  The matrix-valued
connection and its transformation law are independent of this basis
normalization.

\end{bankssolutions}
